\section{Introduction}
\label{sec:introduction}

The seats-votes curve is the foundational concept in the political science literature
on partisan fairness in redistricting. Before the Efficiency Gap, before Declination,
before the Mean-Median Difference---the seats-votes relationship was the starting
point for every analysis of whether redistricting produced systematically asymmetric
partisan outcomes. The curve maps statewide vote share to seat share, and its shape
encodes everything that matters for partisan fairness: a symmetric curve (the ideal)
treats both parties identically at every vote level; an asymmetric curve (the signature
of gerrymandering) favors one party across the full range of plausible election outcomes.

The scalar metrics in L.1--L.4 are all summaries of the seats-votes curve. The
Efficiency Gap is the area between the observed $S(v)$ curve and the diagonal $S(v) = v$.
Partisan Bias is the value of $S(0.50) - 0.50$. Mean-Median Difference is a summary
of the vote-share distribution that affects the shape of the curve near the median.
Declination is a summary of the angular structure of the curve's inflection near 50\%.
None of these scalar summaries captures the full picture that the curve provides.

This paper argues that the seats-votes curve---supplemented by its scalar summary,
Responsiveness ($R = dS/dv|_{v=0.50}$)---is the preferred exhibit for courtroom
presentations of partisan fairness evidence. The curve is the sufficient statistic:
all other L-series metrics can be computed from it, which means reporting the curve
is at least as informative as reporting any combination of the scalar metrics.
The graphical format communicates directly to judicial audiences. And the comparison
of \textsc{bisect} and enacted curves---two step-functions on the same axes,
one passing symmetrically through $(0.50, 0.50)$ with slope $\approx 2.0$, the
other passing below $(0.50, 0.50)$ with slope $\approx 1.3$---provides the most
complete visual evidence of partisan asymmetry available.

Section~\ref{sec:definition} develops the mathematical definition and proves the
sufficiency claim analytically. Section~\ref{sec:behavior} explains why compactness
produces high responsiveness and why enacted safe-seat maps have suppressed responsiveness.
Section~\ref{sec:empirical} presents the empirical curves for NC with responsiveness
and bias statistics. Section~\ref{sec:legal} documents the curve's legal history.
Section~\ref{sec:conclusion} concludes with courtroom guidance.
